% ============================================================================
% NEW MATERIAL, October 6, 2026 (agent draft; not in JMP_DS_draft).
% Adapted from sections/04_quantification.tex (Sept 23, organized after
% Boar, Gorea and Midrigan), which is left unchanged. Numbers are the base
% point promoted Oct 6: output/model/production/2007 (Oct 6 refit, Mac chain 7,
% case 0129_nm; one-birth Estate A + net-estate floor, balanced property-tax
% rebate, PSID levels entry; loss 15.0712 on the 14-row contract; provisional,
% fresh native postcheck passed). Read with
%   python3 code/model/tools/read_results.py fit base_2007
% Conversions traced in tmp/rebate_entry_impl_20261006/closure_root/code/model
% (inputs.py:76 beta; e5f_social_security.py:41 income; parameters.py:1237
% child departure; tools/e5f_initial_fertility_observer.py age windows;
% engine/distribution.py:2435 bequest annualization).
% ============================================================================
\providecommand{\mocktodo}[1]{\textcolor{red}{[TODO: #1]}}

\section{Quantification}
\label{sec:quant-oct06}

I calibrate a stationary economy representing 2007, before the subsequent
decline in fertility. I use observations from the surrounding years to
discipline fertility, housing, and household wealth. Fertility moments come
from the 2004 and 2006 CPS and the 2003--2006 natality records, housing moments
from the 2005--2006 ACS, the 2007 AHS and the PSID, and wealth from the 2005
and 2007 PSID and the 2007 SCF. This steady state is the starting point for the
historical transition.

I organize the parameters into two groups. The first is set outside the
estimation, from external evidence or direct measurement. The second is
estimated by matching model moments to their empirical counterparts. I first
describe the data behind the targets. Before describing either group of
parameters, I then explain how annual data map into the model's
four-year period, because almost every parameter and moment passes through
that mapping.

\subsection{Data}
\label{sec:quant-data}

The targets combine five data sets, each chosen for the object it measures
best. Table~\ref{tab:quant-fit-oct06} lists every moment with its source.

\paragraph{Fertility.} The June Fertility Supplement of the CPS records the
number of children ever born to each woman. I pool the 2004 and 2006
supplements and measure childlessness and the share of mothers with exactly
one child among women aged 40 to 44, and the mean number of children ever born
among women aged 25. Children ever born are capped at three, the largest
number the model distinguishes. The timing of first births comes from the
NCHS natality files, which record every birth with the mother's age and the
birth order; I pool first births in 2003 to 2006.

\paragraph{Housing.} Mean rooms come from the 2007 AHS, for household heads
aged 18 to 85, without top-coding. Ownership among household heads aged 30 to
55 and the ownership gap between recent parents, whose oldest child is under
four, and heads without children come from the 2005--2006 ACS.
\mocktodo{state the dwelling-structure restriction used in both ACS ownership
moments (UNITSSTR codes 3--10 in the builder) and whether the model imposes it.}
The room response to the first birth is the PSID event-study estimate of
Section~\ref{sec:empirical-evidence} in the sample that fixes household
position before the birth: 1.465 rooms three to four years after the birth
(standard error 0.050). I use this sample rather than the household sample of
Figure~\ref{fig:emp-rooms-path}, which gives 1.03, because the model follows
the same households before and after the birth.

\paragraph{Wealth.} Aggregate net worth relative to aggregate annual gross
labor earnings comes from the 2005 and 2007 PSID: net worth of all household
heads aged 18 to 85, excluding business and farm equity and other real
estate, over gross labor earnings of heads aged 18 to 65. The bequest flow is
the annual flow of estates directed to children relative to aggregate net
worth in the 2007 Survey of Consumer Finances.
% [DATA] builders: CPS code/data/cps_fertility/; NCHS
% code/data/nchs_natality_timing/build_first_birth_timing.R; AHS
% code/data/ahs_supply_snapshot/build_ahs_2007_room_target.py; ACS
% code/empirical/housing/build_initial_housing_profile_diagnostic.py; PSID rooms
% code/data/psid_followup_mar2026/sa_rooms_first_birth_v2.do (A2h); SCF
% code/data/scf/build_bequest_flow_2007.py. Target fingerprint c7a3d185...

\subsection{From annual data to four-year periods}
\label{sec:quant-period}

A model period lasts four years. The unit of account is mean annual gross
labor earnings of working-age households, so a model value of one corresponds
to \$82,882 in the 2005 and 2007 PSID (2022 dollars).
% [DATA] metadata.json normalizers: psid_annual_gross_labor_age18_65 82881.95;
% model_annual_gross_labor_before_retirement 1.0522.
The mapping follows three rules.

\paragraph{Stocks are measured in annual units; flows are four years long.}
Wealth, house prices and the bequest shifter $\theta_1$ are stocks and are not
rescaled. Income, consumption, rent and every other flow accrue over the whole
period, so a flow in the model is four times its annual counterpart. For
example, an entrant at age 18 with average productivity earns
$4\times(1-0.080)\times0.721=2.65$ per period after the payroll tax, which is
0.72 times mean annual gross earnings.
% [MODEL] e5f_social_security.py:41: income = period_years*(1-tax)*wage*profile;
% first profile cell 0.72055. Annual gross recovered as income/4/(1-tau).
A moment that relates a stock to a flow is computed against the annual flow,
as in the data: the wealth target divides net worth by annual gross earnings,
and the bequest target divides the estate flow over a period by four before
dividing it by aggregate wealth.

\paragraph{Rates compound; hazards come from durations.}
Interest, discounting and depreciation compound over the four years. The gross
return is $R=1.02^4=1.082$, the period discount factor is $\beta=\beta_{\text{ann}}^4$,
and the period depreciation rate is $1-(1-\delta_{\text{ann}})^4$. With
$\beta_{\text{ann}}=0.940$, the period discount factor is $0.780$. The property
tax is the exception: it is a tax on value levied each year, and I set the
period rate to four times the annual rate.%
\footnote{Compounding instead would give a period rate of 4.17 rather than
4.24 percent of value.}
\mocktodo{confirm the linear property-tax convention; the code calls it
explicit and not interchangeable with compounding.}
Probabilities of events are period probabilities. Children leave home after
18 years on average, which I implement as a constant departure probability of
$\lambda=1/4.5=2/9$ per period for each child. Survival probabilities are
four-year probabilities between successive age cells. The probability $\pi_j$
that an attempt to have a child succeeds is the probability of a conception
within four years at that age.
% [MODEL] parameters.py:68,76,1237 (A_m = 18, stage duration 4.5 periods);
% fecundity pi = 1 - 0.02 exp(0.134 (age-18)); docs/model/
% fecundity_schedule_note_20260720.tex fits within-four-year conception
% probabilities (91/84/64% at 30/35/40).
\mocktodo{cite the source of the four-year conception probabilities named in
docs/model/fecundity\_schedule\_note\_20260720.tex, and the life table behind
the survival probabilities.}
The earnings process is estimated directly at the four-year frequency rather
than converted from an annual process. I take non-overlapping four-year means
of household gross labor earnings in the annual PSID waves of 1984--1997,
remove age and year effects, and estimate $\rho=\gamma_1/\gamma_0=0.735$ from
the first two autocovariances of their logarithm. The implied annual
persistence is $0.735^{1/4}=0.926$.
% [DATA] single_process_external_estimate.json: gamma0 0.5085, gamma1 0.3735,
% sigma_eta 0.484, bootstrap rho [0.692, 0.769]; EARNINDRRC.

\paragraph{Age-specific moments use the overlap of data ages with model cells.}
A model household of age $j$ occupies the four-year cell $[a_j,a_j+4)$. A data
statistic for an age window weights each cell by the fraction of the cell that
falls in the window. Within a cell, a household's number of children is
interpolated linearly between its value at the start and at the end of the
cell, which is exact if births are spread uniformly over the four years. Two
examples show how this works. Women aged 40 to 44 overlap the cell
$[38,42)$ for half its length and the cell $[42,46)$ for three-quarters of its
length, and in each cell the stock of children is read at the midpoint of the
overlap. Children ever born at age 25 are read from the cell $[22,26)$,
seven-eighths of the way from its start to its end. For the mean age at first
birth, I date a first birth in the model at the midpoint of its cell (20, 24,
\ldots, 44) and group the natality data into the same four-year cells, so that
both sides use the same representative ages.
% [MODEL] tools/e5f_initial_fertility_observer.py:154 (midpoints), 163-182
% (40-44 overlap weights 0.5 and 0.75; post shares 0.75 and 0.375),
% 185-218 (age 25: post share (25.5-22)/4 = 0.875).

The room response to the first birth is the one moment for which the period
length bounds how closely the model can mimic the data. In the model, I
compare households that have a first birth with otherwise identical
households that do not, one period, or four years, after the birth. The data
compare years three and four after the birth with years two and three before
it. The model's contrast is therefore a four-year response with no
pre-period, while the data's spans about six years.
\mocktodo{decide whether to state the six-year versus four-year difference this
plainly; the measurement review (docs/model/e5f\_target\_measurement\_review\_20260927.md)
calls the model observer an approximation, not the same estimator.}

\subsection{Parameters set outside the estimation}

Table~\ref{tab:quant-assigned-oct06} summarizes the parameters assigned
outside the estimation.

\paragraph{Demographics and preferences.}
Households enter at age 18, retire at 66, and live at most through age 85.
Survival is certain before retirement. Children remain at home for 18 years
on average. Births are possible through age 45, with the success probability
of an attempt falling from 0.98 at 18 to 0.50 at 42. I set risk aversion to
$\sigma=2$ and use the power equivalence scale with child weight $\omega=0.7$
and exponent $\nu=0.7$. The consumption weight $\alpha=0.733$ is the
expenditure share of nonhousing consumption among childless renters in the
2019--2023 CEX. A birth into the state of three or more children counts as
3.60 children when computing completed fertility: the mean number of
children, capped at five, among women aged 40 to 44 with three or more in the
2024 CPS.
% [MODEL INPUT] tfr_top_bin_weight 3.602359 (provisional; author review).

\paragraph{Housing finance and transaction costs.}
The annual real return on the bond is 2 percent, and selling a home costs 6
percent of its value \citep{greaneyDynamicUrbanEconomics2025}. Annual
depreciation is 1.416 percent and the annual property tax is 1.060 percent of
value. Mortgage debt cannot exceed $\phi=0.80$ of the value of the home, so an
owner must hold at least 20 percent equity at the end of the period in which
it buys. There is no payment-to-income limit.
\mocktodo{name the sources of the depreciation and property-tax rates (recorded
as "adopted national inputs").}

\paragraph{Income and pensions.}
The deterministic age profile of earnings is normalized to average one over
working ages. Productivity follows the four-year process described above, with
$\rho=0.735$ and innovation standard deviation $0.484$, approximated with nine
states. Entrants draw productivity from its stationary distribution. Every
retiree receives the same pension. I set the payroll tax so that the pension
equals 22.9 percent of mean annual gross earnings, the ratio of pension
income to gross earnings in the CPS ASEC; the stationary pension budget then
implies a payroll tax of 8.0 percent.
% [MODEL INPUT] pension_period 0.91778 = 4 x 0.22945; tau_pay 0.080281.
\mocktodo{confirm the pension/earnings ratio's CPS ASEC year and definition.}

\paragraph{Housing products and supply.}
Owners choose among homes of $\{2,4,6,8,10\}$ rooms. Renters choose any size
up to $\bar h^R=6$ rooms. The supply of housing has an elasticity of
$\zeta=0.63$ with respect to rent. The supply scale $H_0$ is not a free
parameter: given the other parameters, it is the value at which the housing
market clears at the price consistent with a constant population.
% [MODEL INPUT] xi_supply 0.63 ("fixed provisional external mapping");
% H0 derived = 6.3117 at N0 = 1.
\mocktodo{cite the source of $\zeta=0.63$ (the September deck cites
Baum-Snow and Han (2024); check before citing).}

\paragraph{Entry wealth and bequests.}
Entrants' wealth follows the PSID rule of Section~\ref{sec:model-oct05}:
net worth of childless renters aged 18 to 24, censored at zero and assigned to
productivity states by income rank. Under Estate A the bequest motive does not
depend on the number of children. I fix the bequest shifter $\theta_1$ at 1
percent of median annual gross earnings, $\theta_1=0.008$.

\begin{table}[H]
\centering
\caption{Parameters set outside the estimation}
\label{tab:quant-assigned-oct06}
\small
\begin{tabularx}{\textwidth}{@{}X l X@{}}
\toprule
Parameter & Value & Evidence or restriction \\
\midrule
Model period & 4 years & Model timing \\
Entry, retirement, maximum ages & 18, 66, 85 & Life-cycle specification \\
Child departure per period $\lambda$ & $2/9$ & Mean stay of 18 years \\
Success of a birth attempt $\pi_j$ & 0.98 (18) to 0.50 (42) & Four-year conception probabilities \\
Risk aversion $\sigma$ & 2 & Maintained \\
Equivalence scale $(\omega,\nu)$ & $(0.7,0.7)$ & Maintained power scale \\
Consumption weight $\alpha$ & 0.733 & CEX 2019--2023, childless renters \\
Annual real return & 2\% & \citet{greaneyDynamicUrbanEconomics2025} \\
Annual depreciation; property tax & 1.416\%; 1.060\% & National inputs \\
Financed share $\phi$ & 0.80 & 20\% equity requirement \\
Selling cost $\varsigma$ & 6\% & \citet{greaneyDynamicUrbanEconomics2025} \\
Pension / mean annual earnings & 0.229 & CPS ASEC; implies payroll tax 8.0\% \\
Persistence $\rho$ (four-year) & 0.735 & PSID 1984--1997 \\
Innovation SD $\sigma_\eta$ (four-year) & 0.484 & PSID 1984--1997 \\
Owner house sizes & 2, 4, 6, 8, 10 rooms & \\
Largest rental unit $\bar h^R$ & 6 rooms & \\
Housing-supply elasticity $\zeta$ & 0.63 & \\
Entry wealth & PSID levels & Childless renters 18--24, censored at zero \\
Bequest shifter $\theta_1$ & 0.008 & 1\% of median annual earnings \\
\bottomrule
\end{tabularx}
\end{table}

\subsection{Internally estimated parameters}

I estimate ten parameters jointly to match ten moments, listed in
Table~\ref{tab:quant-fit-oct06}. In addition, the stationary population
requires mean completed fertility of 2.1, which the house price delivers in
equilibrium. Although all parameters are identified jointly, I organize the
discussion by economic block and describe which moments are most informative
about each.

\paragraph{Saving.}
The annual discount factor governs how much households carry across ages.
Aggregate net worth relative to aggregate annual gross earnings is the most
informative moment. The target is a ratio of totals, not an average of
household ratios.

\paragraph{Housing demand and tenure.}
The room response around the first birth is the most informative moment for
the parent floor $h_P$: the larger the floor, the more rooms a household must
add when its first child arrives. Mean rooms disciplines housing demand more
broadly. The ownership rate and the ownership gap between recent parents and
childless households are most informative about the ownership premium $\chi$
and the dispersion $\kappa_H$ of tenure taste shocks.

\paragraph{Fertility.}
Childlessness is most informative about the first-birth cost $\xi$. The mean
age at first birth and the number of children by age 25 discipline the timing
of first births, and with it the dispersion $\kappa_1$ of first-birth taste
shocks. The share of mothers with exactly one child is most informative about
the dispersion $\kappa_C$ for later births and the curvature $\gamma$ of the
child benefit. The level of the child benefit $\psi$ moves all fertility
moments together.

\paragraph{Bequests.}
With $\theta_1$ fixed, the bequest flow relative to aggregate wealth is most
informative about $\theta_0$.

\begin{table}[H]
\centering
\caption{Internally estimated parameters}
\label{tab:quant-parameters-oct06}
\small
\begin{tabularx}{\textwidth}{@{}l X r X@{}}
\toprule
Parameter & Interpretation & Value & Most informative moments \\
\midrule
$\beta_{\text{ann}}$ & Annual discount factor & 0.940 & Wealth / earnings \\
$\xi$ & First-birth utility cost & 0.309 & Childlessness \\
$\kappa_1$ & First-birth taste dispersion & 0.121 & First-birth age; children by 25 \\
$\kappa_C$ & Later-birth taste dispersion & 0.356 & Exactly one child \\
$\psi$ & Child benefit, level & 0.157 & All fertility moments \\
$\gamma$ & Child benefit, curvature & 0.098 & Exactly one child \\
$h_P$ & Parent floor (rooms) & 2.585 & First-birth room response; mean rooms \\
$\chi$ & Ownership premium & 1.081 & Ownership; recent-parent gap \\
$\kappa_H$ & Tenure taste dispersion & 0.018 & Ownership; recent-parent gap \\
$\theta_0$ & Bequest strength & 0.212 & Bequests / wealth \\
\addlinespace
$H_0$ & Supply scale (derived) & 6.312 & Market clearing at constant population \\
\bottomrule
\end{tabularx}
\par\smallskip
\begin{minipage}{0.95\linewidth}
\footnotesize\textit{Notes:} Search bounds: $\beta_{\text{ann}}\in[0.93,0.99]$,
$h_P\in[0.1,2.6]$, $\kappa_1,\kappa_C\in[0.02,50]$.
\mocktodo{$h_P$ is within 0.015 of its upper bound and the search flags
$\kappa_1$ and $\kappa_C$ as near their bounds; decide how to report this.
$\beta_{\text{ann}}=0.9397$ is below the 0.94 floor recorded as an external
restriction in July; confirm which restriction applies.}
% [MODEL] output/model/production/2007/parameters_estate_a.csv.
\end{minipage}
\end{table}

\subsection{Target moments and fit}

Table~\ref{tab:quant-fit-oct06} reports the targets and the model's values.
The model matches childlessness, the share with exactly one child, the timing
of first births, ownership, the recent-parent ownership gap and the wealth and
bequest ratios closely. It misses two moments. Households in the model have
0.55 children by age 25 against 0.81 in the data, so first births start too
late even though their mean age is matched. And the model's family adds 1.26
rooms around the first birth against 1.47 in the data, with the floor $h_P$
near its upper bound.

\begin{table}[t]
\centering
\caption{Targets and model moments, 2007 steady state}
\label{tab:quant-fit-oct06}
\small
\begin{tabularx}{\textwidth}{@{}X l r r@{}}
\toprule
Moment & Data & Target & Model \\
\midrule
\multicolumn{4}{@{}l}{\emph{Fertility}}\\
Childlessness, women 40--44 (\%) & CPS 2004, 2006 & 19.83 & 20.20 \\
Exactly one child among mothers 40--44 (\%) & CPS 2004, 2006 & 21.37 & 21.76 \\
Mean age at first birth (years) & NCHS 2003--2006 & 25.98 & 25.99 \\
Children ever born at age 25 & CPS 2004, 2006 & 0.810 & 0.551 \\
First births at age 30+ (\%)$^{\dagger}$ & NCHS 2003--2006 & 24.93 & 23.88 \\
\addlinespace
\multicolumn{4}{@{}l}{\emph{Housing}}\\
Mean rooms, heads 18--85 & AHS 2007 & 5.729 & 5.816 \\
Ownership, heads 30--55 (\%) & ACS 2005--2006 & 67.63 & 66.68 \\
Recent-parent ownership gap (pp) & ACS 2005--2006 & 12.76 & 12.72 \\
Rooms added around the first birth & PSID 1968--2019 & 1.465 & 1.260 \\
Rooms, 3+ minus 1--2 children$^{\dagger}$ & ACS 2005--2006 & 0.385 & 0.295 \\
\addlinespace
\multicolumn{4}{@{}l}{\emph{Wealth}}\\
Net worth / annual gross earnings & PSID 2005, 2007 & 4.458 & 4.576 \\
Bequests / wealth (\%) & SCF 2007 & 0.729 & 0.702 \\
p90/p50 wealth-to-income, ages 76--84$^{\dagger}$ & PSID 2003--2005 & 3.516 & 3.307 \\
\addlinespace
Mean completed fertility (equilibrium condition) & --- & 2.10 & 2.10 \\
\bottomrule
\end{tabularx}
\par\smallskip
\begin{minipage}{0.95\linewidth}
\footnotesize\textit{Notes:} $^{\dagger}$ Not targeted. Samples and definitions in
Section~\ref{sec:quant-data}.
\end{minipage}
\end{table}
% Full fit (base_2007, loss 15.0712): moment | target | model | gap | weight | loss
% cps_childlessness .198279 .202033 +.003755 35532.3 0.5009
% cps_exactly_one .213655 .217574 +.003919 26952.8 0.4139
% nchs_mean_age 25.97626 25.98889 +.012624 139.828 0.0223
% nchs_share30 .249278 .238803 -.010475 0 0 (validation)
% wealth_earnings 4.458387 4.576325 +.117938 7.5951 0.1056
% bequest_wealth .0072910 .0070194 -.000272 5165289.3 0.3811
% old_dispersion 3.515935 3.306923 -.209012 0 0 (validation)
% mean_rooms 5.729434 5.816323 +.086888 128.021 0.9665
% ownership_30_55 .676260 .666805 -.009455 2339.36 0.2091
% first_birth_rooms 1.465 1.259836 -.205164 137.565 5.7904
% family_rooms .385100 .295259 -.089841 0 0 (validation)
% recent_parent_ownership .127608 .127185 -.000423 27055.8 0.0048
% early_fertility .809528 .551139 -.258389 100 6.6765
% initial_normalization 2.1 2.1000001 (normalization, unscored)
